\newcommand{\figuraRecipienteRotante}{%
\begin{figure}[H]
    \centering
    \begin{tikzpicture}[>=Latex,font=\small,x=1.15cm,y=0.95cm,
                        line cap=round,line join=round]
        \draw[very thick] (-4,4.8)--(-4,0)--(4,0)--(4,4.8);
        \fill[gray!28] (-3.92,0.08)--(3.92,0.08)--
            plot[domain=3.92:-3.92,samples=120] (\x,{2.15+0.13*\x*\x})--cycle;
        \draw[very thick,blue!70!black,domain=-3.92:3.92,samples=120]
            plot (\x,{2.15+0.13*\x*\x});
        \draw[densely dashed] (0,0)--(0,5.45);
        \draw[->,ultra thick,purple!75!black] (0,5.05)
            arc[start angle=90,end angle=-235,radius=0.62];
        \node at (0.95,5.05) {$\Omega$};
        \draw[<->] (-0.45,0)--(-0.45,2.15) node[midway,left] {$z_0$};
        \fill (2.35,3.0) circle (0.10)
            node[above right] {$(r,z)$};
        \draw[->,thin] (0,0.35)--(2.35,0.35) node[midway,below] {$r$};
        \draw[->,thin] (0.35,0)--(0.35,3.0) node[midway,right] {$z$};
    \end{tikzpicture}
    \caption{Corte meridional de un líquido en rotación de cuerpo rígido.
    La superficie libre es parabólica, con altura mínima $z_0$ sobre el eje
    y altura creciente con el radio cilíndrico $r$.}
    \label{fig:recipiente-rotante}
\end{figure}%
}